\begin{result}{D.1}{Known}{The two manuscript assertions have different targets}
In cote 149, scan 56, Conjecture A classifies full profinite sections
of a hyperbolic curve over a field finitely generated over \(\mathbb Q\)
as ordinary rational-point sections or sections in a rational cuspidal
packet. The margin additionally predicts the indicated geometric
fixed-subgroups. In scans 126--127, assertion a) concerns valuation
decomposition groups; assertion b) is the lifting assertion that a
section over an elementary hyperbolic model lifts to the absolute
Galois group of its function field. These are not assertions about the
pronilpotent quotient.
\end{result}
\begin{proof}[Verification of scope]
The supplied second transcription of scan 56 reads
``soit d\'efinie par un pt de \(U(K)\), soit par une section d'une
extension (2)''. The preceding scans 54--55 define the cusp packets.
The text of scan 127 reads
``Alors tte section \ldots\ se rel\`eve en une section'' of the
function-field group. Thus b) is specifically a lifting statement,
while a) supplies the proposed valuation description.
We have read these supplied transcriptions \citep{RootManuscriptTranscriptions},
not the manuscript scans;
uncertain readings in them remain uncertain. No new claim about the
accuracy of the transcription is made here.
\end{proof}

\begin{result}{D.2}{Proved}{Conditional consequences for Conjecture A and birational lifting}
If A2 has an affirmative answer, Conjecture A holds for every full
section of the tripod over \(\mathbb Q\) whose restrictions at all
finite primes are ordinary. Each such section also has a birational
lift, as requested by assertion b) of cote 149, scan 127. This
implication does not settle sections with cuspidal local restrictions,
the additional fixed-subgroup assertions in scan 56, or the main
pronilpotent finite-support question itself.
\end{result}
\begin{proof}
A.5 identifies the full section with \(s_t\) for a rational point
\(t\in U(\mathbb Q)\). To construct a birational lift, let \(z\)
be a uniformizer at that point in \(\mathbb Q(U)\), and choose an
extension of its discrete valuation to an algebraic closure. Its
decomposition group has the exact sequence
\[
 1\longrightarrow\widehat{\mathbb Z}(1)
 \longrightarrow D_t\longrightarrow G_{\mathbb Q}\longrightarrow1.
\]
It splits: in equal characteristic zero, choose compatible roots of
\(z\) and let \(G_{\mathbb Q}\) act on the coefficient field while
fixing these roots. This defines continuous compatible actions on
all finite extensions of the resulting tame local field, hence a
continuous section. Equivalently, use the coefficientwise action on
the algebraic Puiseux extension of \(\overline{\mathbb Q}((z))\).
The inclusion of the decomposition group into
\(G_{\mathbb Q(U)}\) followed by the quotient to \(E_U\) sends
this section to \(s_t\); inertia at the interior point is killed
because the covers of \(U\) are unramified there.
Choices of prolongation alter the construction by geometric
conjugacy. This proves the required existence of a birational lift
of the section class. No assertion about all other full sections is
made by this construction.
\end{proof}

\begin{result}{D.3}{Open}{Unconditional manuscript outcome of this run}
The antecedent in D.2 has not been established. We therefore obtain no
unconditional proof or counterexample to Conjecture A in scan 56 or to
the universal birational lifting assertion b) in scan 127. The precise
unconditional improvement is A.4: for full locally ordinary sections,
the rational-abelian-pair case already follows from an established
theorem and need not remain a second conjectural gate.
\end{result}
